\section{Introduction}
\label{sec:heldout-context}
A software engineering agent must understand a specification, use the relevant code and documentation, construct a solution, and test it. These functions can be distributed among a planner, worker and critic. If separate model calls receive the same material, the context must be transmitted repeatedly. We investigate whether these functions can be retained with fewer calls while leaving correctness checks to external tests.

Hybrid-INoT addresses this question by combining internal role organization with an external verification layer. Planning, implementation and critique are prescribed in a shared model call; an external process evaluates the resulting program. The design draws on INoT's use of internal dialogue~\cite{sun2025inot}, but the role-labelled configuration examined here is distinct from the original PromptCode algorithm. It separates program generation from correctness checking. Long shared context motivates the resource model; it does not by itself guarantee that a single call will preserve quality.

The original Hybrid-INoT implementation also included context selection and selective reruns. The preliminary comparison of complete configurations changed these mechanisms together with call topology, so it cannot identify the contribution of internal roles. The present evaluation isolates the architecture's single-pass core: full supplied context, one final candidate, prescribed planning/implementation/review and subsequent native testing. Compression and verification-driven retries are disabled. Evaluating the components separately tests these architectural choices within the original engineering problem.

The core has two distinct empirical questions. Does assigning role names add value when the stage operations stay fixed? Does performing the operations in one response improve the quality--resource tradeoff relative to distributing them across three calls? A direct solver tests whether the staged organization is useful at all. A separate INoT adaptation examines the related algorithmic alternative. Each condition receives the same supplied task and context and undergoes the same subsequent evaluation. The design never treats a model's self-critique as the measured verification result.

The contributions are:
\begin{enumerate}
\item An operational account of the Hybrid-INoT single-pass core, its state and its external evaluation boundary, with a direct mapping to every implemented control.
\item A conditional shared-context token balance explaining when fewer calls can reduce resources, without assuming equal output lengths or asserting a measured long-context crossover.
\item A compression-free $2\times2$ component study and direct comparator on 200 reserved random BigCodeBench tasks: 996 observed programs, original tests, four predeclared paired quality contrasts and explicit missingness.
\item Complementary evidence from 400 development programs, 199 programs from an INoT adaptation and a one-issue software-repair demonstration, with complete function-generation outcomes and experimental instructions in the appendices.
\end{enumerate}
Role ablation and interaction-count studies have close predecessors. The contribution here is the specified architectural evaluation and its auditable evidence, not the invention of role collaboration. The findings constrain the design: one-call roles have a lower observed resource mean, but quality non-inferiority is not established; three-call organization requires more resources without a detected quality improvement in the four planned tests.

\section{Related work}
\label{sec:heldout-related}
\subsection{Closest methodological antecedent}
Self-collaboration Code Generation via ChatGPT~\cite{dong2023selfcollaboration} already ablates role instructions and maximum interactions in Sections~5.2--5.3. Removing roles or varying interaction is therefore not a novelty claim of the present study. Table~\ref{tab:closest-method} distinguishes the implemented interventions.

\begin{table}[htbp]\centering\small
\caption{Method comparison with the closest antecedent~\cite{dong2023selfcollaboration}, Tables~3--4 and Appendix~B.1. Interaction rounds and fresh model calls are different quantities.}\label{tab:closest-method}
\begin{tabular}{p{0.20\textwidth}p{0.31\textwidth}p{0.36\textwidth}}\toprule
Aspect & Self-collaboration paper & Component experiment\\\midrule
Labels & Role versus role-excised instructions & Identical operation sentences; only the stage prefixes change\\
Interaction & Feedback loop; maximum interaction varied & One or three fresh turns; fixed sequence, full exposed history\\
Design & Separate role and interaction ablations & Joint label-by-call matrix on every assigned task\\
Tasks & HumanEval, MBPP and extended tests & Reserved BigCodeBench sample; pre-generation reference and incorrect controls\\\bottomrule
\end{tabular}\end{table}

The joint design estimates the interaction of role labels and call topology under a fixed generation interface. The evaluated programs, native outcomes and token components make its quality and resource comparisons inspectable. This extends the empirical comparison without claiming a new principle of role collaboration.

Self-collaboration is the closest external method because its role ablation addresses the same attribution question. Our external comparison uses the pinned 2024 author implementation, whose generated-test execution differs from the simulated testing described in the 2023 paper (Section~\ref{sec:scc-feasibility}). The frozen comparison assigns SCC, fresh SR and fresh SN to the same 1,000 tasks with three repetitions; it contains no role-excised SCC arm. It therefore estimates differences between complete workflows, including feedback and stopping, while SR--SN and MR--MN isolate the specified labels within our component design. Two development checks establish execution feasibility, but the main comparison remains unfinished. D provides the minimal comparator. INoT* addresses a separate internal-dialogue algorithm question.

\subsection{Internal dialogue and repository methods}

INoT describes compact PromptCode instructions for internal dialogue and reasoning~\cite{sun2025inot}. This motivates evaluating a disclosed algorithm adaptation, while a separate minimal label intervention addresses whether named roles themselves change results. The two comparisons answer different questions: the INoT instruction changes the reasoning procedure, whereas the factorial label contrast preserves the operation sentences. The original Hybrid-INoT comparison additionally changed compression and reruns; its results are described in the historical appendix and are not pooled with the new experiment.

Single- versus multi-agent work under matched thinking-token budgets emphasizes that additional calls and additional computation are distinct interventions~\cite{tran2026budget}. Here, actual completion length is not held to a matched hard budget. We therefore report token components and list-price valuations alongside quality. These measurements describe the implemented protocols; they do not identify an optimal allocation at a fixed inference budget.

Agentless uses a structured localization--repair--validation workflow~\cite{xia2024agentless}; SWE-agent couples the model to a repository interface~\cite{yang2024sweagent}. Both add retrieval, tools, and environment interaction beyond fixed-context function generation, so their published scores are not directly comparable with this paired matrix. BigCodeBench~\cite{zhuo2024bigcodebench} supplies diverse library-dependent tasks with executable tests, whereas SWE-bench~\cite{jimenez2023swebench,swebenchdocs2026} evaluates repository changes against issue-specific tests. Native evaluation strengthens outcome validity, but passing a reference control does not resolve every disagreement between task prose and tests.

Feedback is a separate design choice from role naming. Self-Refine alternates model-generated feedback and revision using the same model~\cite{madaan2023selfrefine}. Reflexion retains verbal reflections between trials; its programming experiments use self-generated executable tests as feedback~\cite{shinn2023reflexion}. Our component conditions receive no execution feedback during generation and follow a fixed sequence, so they do not reproduce either adaptive procedure. The SCC author-code comparison adds generated-test feedback within its complete workflow. In our experiments, benchmark tests and reference programs remain outside the generation interface; internal SCC execution never determines benchmark success. These distinctions make feedback access and stopping part of the external method comparison, rather than evidence for a role-label effect. Self-Refine and Reflexion provide methodological context; neither has been run as an experimental arm here.

Persona studies motivate a narrower claim than ``roles improve reasoning.'' Zheng et al. find that adding system personas does not reliably improve their factual-question evaluations~\cite{zheng2024personas}. Luz de Araujo et al. distinguish performance, robustness to irrelevant persona attributes, and fidelity to the assigned expertise; their results vary across models and tasks~\cite{araujo2025personas}. These are not tests of our coding prompts. They motivate changing labels while retaining the requested operations, and reporting the result for the exact labels used rather than for all possible personas.

Multi-agent evidence also depends on its comparator. Choi et al. attribute much of the gain in their NLP debate experiments to independent voting and show why communication must be separated from aggregation~\cite{choi2025debate}. Conversely, Wunderlich et al. report configurations where debate or mixture-of-agents improve the compute--accuracy tradeoff relative to self-consistency on MMLU-Pro and BBH~\cite{wunderlich2026compute}. Together these findings support conditional evaluation, not a universal claim for or against multiple agents. Our neutral serial pipeline is not voting, a mixture of independent models, or a matched-budget ensemble. This study provides a within-model factorial test on executable code, preserving traces and accounting for available resource counters; it does not estimate general superiority over those alternatives.

